\documentclass[12pt, a4paper]{article}

% Packages
\usepackage{amsmath, amssymb, amsthm}
\usepackage{physics}
\usepackage{geometry}
\usepackage{graphicx}
\usepackage{booktabs}
\usepackage{array}
\usepackage{hyperref}
\usepackage{setspace}
\usepackage{enumitem}
\geometry{margin=2.5cm}
\setstretch{1.2}

\begin{document}

\section*{Problem}

In the framework of quantum chromodynamics (QCD), monopole solutions are
important topological soliton configurations. Here we consider the non-Abelian
generalization of the Dirac monopole, namely the Wu--Yang monopole in an
$SU(2)$ Yang--Mills theory on four-dimensional Minkowski spacetime.

The Yang--Mills Lagrangian is
\[
L = -\frac{1}{4} \int d^4 x \sum_{a=1}^{3} F^a_{\mu\nu} F^{a,\mu\nu}.
\]

The Minkowski metric is chosen as
\[
\eta_{\mu\nu} = \eta^{\mu\nu}
=
\begin{pmatrix}
1 & 0 & 0 & 0 \\
0 & -1 & 0 & 0 \\
0 & 0 & -1 & 0 \\
0 & 0 & 0 & -1
\end{pmatrix}.
\]

We introduce the generalized Wu--Yang monopole gauge field
\[
A_0^a = 0,
\qquad
A_i^a = \sum_{b=1}^{3} \epsilon_{aib}\,\frac{x^b}{g\,r^2}\qty[1 + \phi(r)],
\]
with
\[
r^2 = (x^1)^2 + (x^2)^2 + (x^3)^2.
\]
Here $\epsilon_{aib}$ is the Levi-Civita symbol in three spatial dimensions
and $g$ is the Yang--Mills coupling constant.

Under the condition $\phi'(0)=0$, derive the second-order differential
equation satisfied by $\phi(r)$ and express $\phi''(r)$ purely in terms of
$\phi(r)$ and $r$.

\section*{Answer}

\[
\phi''(r) = \frac{\phi(r)\qty[\phi^2(r)-1]}{r^2}.
\]

\section*{Derivation}

Since there is no electric field and the configuration is static, only the
spatial components of the gauge field are nonzero. Writing out the gauge
potential explicitly gives the following nonzero components:

\[
\begin{aligned}
&\text{For } a=1: &&
A_1^1 = 0, \qquad
A_2^1 = \frac{x^3}{g\,r^2}\qty[1+\phi(r)], \qquad
A_3^1 = -\frac{x^2}{g\,r^2}\qty[1+\phi(r)], \\[4pt]
&\text{For } a=2: &&
A_1^2 = -\frac{x^3}{g\,r^2}\qty[1+\phi(r)], \qquad
A_2^2 = 0, \qquad
A_3^2 = \frac{x^1}{g\,r^2}\qty[1+\phi(r)], \\[4pt]
&\text{For } a=3: &&
A_1^3 = \frac{x^2}{g\,r^2}\qty[1+\phi(r)], \qquad
A_2^3 = -\frac{x^1}{g\,r^2}\qty[1+\phi(r)], \qquad
A_3^3 = 0.
\end{aligned}
\]

% [Please verify: the PDF writes the monopole ansatz with Levi-Civita symbols;
% the explicit component signs above are the convention consistent with the
% subsequent field-strength identities and the final equation of motion.]

The Yang--Mills field strength is defined by
\[
F_{\mu\nu}^a
=
\partial_\mu A_\nu^a
- \partial_\nu A_\mu^a
- g\,\epsilon_{abc} A_\mu^b A_\nu^c.
\]

We first compute the derivative of the scalar factor:
\[
\partial_i r = \pdv{r}{x^i} = \frac{x^i}{r},
\]
and therefore
\[
\partial_i\qty(\frac{1+\phi(r)}{r^2})
=
\frac{x^i}{r}\,\dv{r}\qty(\frac{1+\phi(r)}{r^2})
=
\frac{x^i}{r^4}\qty[r\phi'(r)-2r(1+\phi(r))]
=
\frac{x^i}{r^4}\qty[r\phi'(r)-2\phi(r)-2].
\]

\subsection*{Color component $a=1$}

For the first color component, we compute:
\[
\begin{aligned}
F_{12}^1
&=
\partial_1 A_2^1 - \partial_2 A_1^1 - g\,\epsilon_{1bc} A_1^b A_2^c \\
&=
\partial_1 A_2^1 - g\,\epsilon_{123} A_1^2 A_2^3 - g\,\epsilon_{132} A_1^3 A_2^2 \\
&=
\partial_1 A_2^1 - g A_1^2 A_2^3 \\
&=
\frac{x^1 x^3}{g\,r^4}\qty[r\phi'(r)-2\phi(r)-2]
- g\qty[-\frac{x^3}{g\,r^2}\qty(1+\phi(r))]\qty[-\frac{x^1}{g\,r^2}\qty(1+\phi(r))] \\
&=
\frac{x^1 x^3}{g\,r^4}\qty[r\phi'(r)-2\phi(r)-2]
- \frac{x^1 x^3}{g\,r^4}\qty[1+\phi(r)]^2 \\
&=
-\frac{x^1 x^3}{g\,r^4}\qty[r\phi'(r)+\phi^2(r)-1]
= -F_{21}^1.
\end{aligned}
\]

\[
\begin{aligned}
F_{23}^1
&=
\partial_2 A_3^1 - \partial_3 A_2^1 - g\,\epsilon_{1bc} A_2^b A_3^c \\
&=
\partial_2 A_3^1 - \partial_3 A_2^1
- g\,\epsilon_{123} A_2^2 A_3^3
- g\,\epsilon_{132} A_2^3 A_3^2 \\
&=
\partial_2 A_3^1 - \partial_3 A_2^1 + g A_2^3 A_3^2 \\
&=
-\partial_2\qty[\frac{x^2}{g\,r^2}\qty(1+\phi(r))]
-\partial_3\qty[\frac{x^3}{g\,r^2}\qty(1+\phi(r))]
+ g\qty[\frac{x^1}{g\,r^2}\qty(1+\phi(r))]\qty[\frac{x^1}{g\,r^2}\qty(1+\phi(r))] \\
&=
-\frac{1}{g\,r^4}\qty[(r^2-(x^2)^2)\qty(1+\phi(r)) + x_2^2\qty(r\phi'(r)-2\phi(r)-2)] \\
&\quad
-\frac{1}{g\,r^4}\qty[(r^2-(x^3)^2)\qty(1+\phi(r)) + x_3^2\qty(r\phi'(r)-2\phi(r)-2)] \\
&\quad
+ \frac{(x^1)^2}{g\,r^4}\qty[1+\phi(r)]^2 \\
&=
\frac{1}{g\,r^4}
\qty[
\qty(r^2-(x^1)^2) r\phi'(r)
+ \qty(\phi^2(r)-1)(x^1)^2
] \\
&=
-F_{32}^1.
\end{aligned}
\]

\[
\begin{aligned}
F_{13}^1
&=
\partial_1 A_3^1 - \partial_3 A_1^1 - g\,\epsilon_{1bc} A_1^b A_3^c \\
&=
\partial_1 A_3^1
- g\,\epsilon_{123} A_1^2 A_3^3
- g\,\epsilon_{132} A_1^3 A_3^2 \\
&=
\partial_1 A_3^1 + g A_1^3 A_3^2 \\
&=
-\frac{x^1 x^2}{g\,r^4}\qty[r\phi'(r)-2\phi(r)-2]
+ g\qty[\frac{x^2}{g\,r^2}\qty(1+\phi(r))]\qty[\frac{x^1}{g\,r^2}\qty(1+\phi(r))] \\
&=
\frac{x^1 x^2}{g\,r^4}\qty[r\phi'(r)+\phi^2(r)-1]
= -F_{31}^1.
\end{aligned}
\]

\subsection*{Color component $a=2$}

Similarly,
\[
\begin{aligned}
F_{12}^2
&=
\partial_1 A_2^2 - \partial_2 A_1^2 - g\,\epsilon_{2bc} A_1^b A_2^c \\
&=
-\partial_2 A_1^2 - g\,\epsilon_{213} A_1^1 A_2^3 - g\,\epsilon_{231} A_1^3 A_2^1 \\
&=
-\partial_2 A_1^2 - g A_1^3 A_2^1 \\
&=
\frac{x^2 x^3}{g\,r^4}\qty[r\phi'(r)-2\phi(r)-2]
- g\qty[\frac{x^2}{g\,r^2}\qty(1+\phi(r))]\qty[\frac{x^3}{g\,r^2}\qty(1+\phi(r))] \\
&=
-\frac{x^2 x^3}{g\,r^4}\qty[r\phi'(r)+\phi^2(r)-1]
= -F_{21}^2.
\end{aligned}
\]

\[
\begin{aligned}
F_{23}^2
&=
\partial_2 A_3^2 - \partial_3 A_2^2 - g\,\epsilon_{2bc} A_2^b A_3^c \\
&=
\partial_2 A_3^2 - g\,\epsilon_{213} A_2^1 A_3^3 - g\,\epsilon_{231} A_2^3 A_3^1 \\
&=
\partial_2 A_3^2 - g A_2^3 A_3^1 \\
&=
\frac{x^1 x^2}{g\,r^4}\qty[r\phi'(r)-2\phi(r)-2]
- g\qty[\frac{x^1}{g\,r^2}\qty(1+\phi(r))]\qty[-\frac{x^2}{g\,r^2}\qty(1+\phi(r))] \\
&=
\frac{x^1 x^2}{g\,r^4}\qty[r\phi'(r)+\phi^2(r)-1]
= -F_{32}^2.
\end{aligned}
\]

\[
\begin{aligned}
F_{13}^2
&=
\partial_1 A_3^2 - \partial_3 A_1^2 - g\,\epsilon_{2bc} A_1^b A_3^c \\
&=
\partial_1 A_3^2 - \partial_3 A_1^2
- g\,\epsilon_{213} A_1^1 A_3^3
- g\,\epsilon_{231} A_1^3 A_3^1 \\
&=
\partial_1 A_3^2 - \partial_3 A_1^2 - g A_1^3 A_3^1 \\
&=
\partial_1\qty[\frac{x^1}{g\,r^2}\qty(1+\phi(r))]
+ \partial_3\qty[\frac{x^3}{g\,r^2}\qty(1+\phi(r))]
- g\qty[\frac{x^2}{g\,r^2}\qty(1+\phi(r))]\qty[-\frac{x^2}{g\,r^2}\qty(1+\phi(r))] \\
&=
\frac{1}{g\,r^4}
\qty[
\qty(r^2-(x^2)^2) r\phi'(r)
+ \qty(\phi^2(r)-1)(x^2)^2
] \\
&=
-F_{31}^2.
\end{aligned}
\]

\subsection*{Color component $a=3$}

Finally,
\[
\begin{aligned}
F_{12}^3
&=
\partial_1 A_2^3 - \partial_2 A_1^3 - g\,\epsilon_{3bc} A_1^b A_2^c \\
&=
-\partial_1 A_2^3 - \partial_2 A_1^3
- g\,\epsilon_{312} A_1^1 A_2^2
- g\,\epsilon_{321} A_1^2 A_2^1 \\
&=
-\partial_1 A_2^3 - \partial_2 A_1^3 + g A_1^2 A_2^1 \\
&=
\frac{1}{g\,r^4}
\qty[
\qty(r^2-(x^3)^2) r\phi'(r)
+ \qty(\phi^2(r)-1)(x^3)^2
] \\
&=
-F_{21}^3.
\end{aligned}
\]

\[
\begin{aligned}
F_{23}^3
&=
\partial_2 A_3^3 - \partial_3 A_2^3 - g\,\epsilon_{3bc} A_2^b A_3^c \\
&=
-\partial_3 A_2^3
- g\,\epsilon_{312} A_2^1 A_3^2
- g\,\epsilon_{321} A_2^2 A_3^1 \\
&=
-\partial_3 A_2^3 - g A_2^1 A_3^2 \\
&=
-\frac{x^1 x^3}{g\,r^4}\qty[r\phi'(r)-2\phi(r)-2]
- g\qty[\frac{x^3}{g\,r^2}\qty(1+\phi(r))]\qty[\frac{x^1}{g\,r^2}\qty(1+\phi(r))] \\
&=
-\frac{x^1 x^3}{g\,r^4}\qty[r\phi'(r)+\phi^2(r)-1]
= -F_{32}^3.
\end{aligned}
\]

\[
\begin{aligned}
F_{13}^3
&=
\partial_1 A_3^3 - \partial_3 A_1^3 - g\,\epsilon_{3bc} A_1^b A_3^c \\
&=
-\partial_3 A_1^3
- g\,\epsilon_{312} A_1^1 A_3^2
- g\,\epsilon_{321} A_1^2 A_3^1 \\
&=
-\partial_3 A_1^3 + g A_1^2 A_3^1 \\
&=
\frac{x^2 x^3}{g\,r^4}\qty[r\phi'(r)-2\phi(r)-2]
- g\qty[-\frac{x^3}{g\,r^2}\qty(1+\phi(r))]\qty[-\frac{x^2}{g\,r^2}\qty(1+\phi(r))] \\
&=
-\frac{x^2 x^3}{g\,r^4}\qty[r\phi'(r)+\phi^2(r)-1]
= -F_{31}^3.
\end{aligned}
\]

\subsection*{Contraction of the field strength}

Since the only nonzero components are spatial,
\[
F^{a,\mu\nu} = \eta^{\mu\rho}\eta^{\nu\sigma} F^a_{\rho\sigma},
\qquad
\mu,\nu = 1,2,3,
\]
and therefore
\[
F^{a,ij} = \eta^{ik}\eta^{jl} F^a_{kl}
= (-\delta^{ik})(-\delta^{jl}) F^a_{kl}
= F^a_{ij}.
\]

Hence the Yang--Mills density becomes
\[
\begin{aligned}
\mathcal{L}_{YM}
&=
\sum_{a=1}^{3} F^a_{\mu\nu} F^{a,\mu\nu} \\
&=
2\qty[(F_{12}^1)^2 + (F_{13}^1)^2 + (F_{23}^1)^2] \\
&\quad
+ 2\qty[(F_{12}^2)^2 + (F_{13}^2)^2 + (F_{23}^2)^2] \\
&\quad
+ 2\qty[(F_{12}^3)^2 + (F_{13}^3)^2 + (F_{23}^3)^2].
\end{aligned}
\]

Denote these three color contributions by
\[
\mathcal{L}_{YM} = \mathcal{L}_1 + \mathcal{L}_2 + \mathcal{L}_3.
\]

For $a=1$ one obtains
\[
\mathcal{L}_1
=
\frac{2}{g^2 r^8}
\qty[
r^2 \phi'^2(r)\qty(r^2-(x^1)^2)
+ (x^1)^2\qty(\phi^2(r)-1)^2
].
\]

By symmetry, the $a=2$ and $a=3$ parts are
\[
\mathcal{L}_2
=
\frac{2}{g^2 r^8}
\qty[
r^2 \phi'^2(r)\qty(r^2-(x^2)^2)
+ (x^2)^2\qty(\phi^2(r)-1)^2
],
\]
\[
\mathcal{L}_3
=
\frac{2}{g^2 r^8}
\qty[
r^2 \phi'^2(r)\qty(r^2-(x^3)^2)
+ (x^3)^2\qty(\phi^2(r)-1)^2
].
\]

Adding them together gives
\[
\begin{aligned}
\mathcal{L}_{YM}
&=
\frac{2}{g^2 r^8}\Bigl[
r^2 \phi'^2(r)\Bigl(
\qty(r^2-(x^1)^2)
+ \qty(r^2-(x^2)^2)
+ \qty(r^2-(x^3)^2)
\Bigr) \\
&\quad
+ \qty((x^1)^2 + (x^2)^2 + (x^3)^2)\qty(\phi^2(r)-1)^2
\Bigr] \\
&=
\frac{2}{g^2 r^8}
\qty{
2 r^4 \phi'^2(r) + r^2 \qty(\phi^2(r)-1)^2
} \\
&=
\frac{4}{g^2 r^4}
\qty[
r^2 \phi'^2(r) + \frac{1}{2}\qty(\phi^2(r)-1)^2
] \\
&=
\frac{4}{g^2 r^2}\phi'^2(r)
+ \frac{2}{g^2 r^4}\qty(\phi^2(r)-1)^2.
\end{aligned}
\]

\subsection*{Reduced radial action}

Using spherical coordinates, the four-dimensional integral becomes
\[
\int d^4 x
=
\int dt \int d^3 x
=
\int dt \int_0^{2\pi} d\varphi \int_0^\pi \sin\theta\,d\theta \int_0^\infty r^2\,dr
=
4\pi \int dt \int_0^\infty r^2\,dr.
\]

Therefore,
\[
\begin{aligned}
L
&=
-\frac{1}{4}\,4\pi \int dt \int_0^\infty
r^2 \qty[
\frac{4}{g^2 r^2}\phi'^2(r)
+ \frac{2}{g^2 r^4}\qty(\phi^2(r)-1)^2
]\,dr \\
&=
-\frac{4\pi}{g^2} \int dt \int_0^\infty
\qty[
\phi'^2(r)
+ \frac{\qty(\phi^2(r)-1)^2}{2r^2}
]\,dr.
\end{aligned}
\]

Define the reduced radial density
\[
\mathcal{L}(\phi,\phi')
=
\phi'^2(r) + \frac{\qty(\phi^2(r)-1)^2}{2r^2}.
\]

The Euler--Lagrange equation is
\[
\pdv{\mathcal{L}}{\phi(r)} - \dv{r}\qty(\pdv{\mathcal{L}}{\phi'(r)}) = 0.
\]

Its two ingredients are
\[
\pdv{\mathcal{L}}{\phi(r)}
=
\frac{2\qty[\phi^2(r)-1]\phi(r)}{r^2},
\qquad
\pdv{\mathcal{L}}{\phi'(r)}
=
2\phi'(r).
\]

Substituting into the Euler--Lagrange equation yields
\[
\frac{2\qty[\phi^2(r)-1]\phi(r)}{r^2} - 2\phi''(r) = 0.
\]

Hence,
\[
\boxed{
\phi''(r) = \frac{\phi(r)\qty[\phi^2(r)-1]}{r^2}
}.
\]

\end{document}
